\lecture{22}
\title{Bayesian Inference II}

\begin{document}

\subsection{Computing the Bayes Estimator and MAP}

\begin{quote}
    \textbf{Question \#1.} Determine how one might compute the Bayes estimator
    $\hat{\theta}^{\text{Bayes}}
    := \int \theta \cdot f(\theta \mid X_1, \dots, X_n)
    \, \mathrm{d}\theta$.
\end{quote}

Here's the first idea:
instead of computing an integral,
just sample $N$ points according to the distribution
$f(\theta \mid X_1, \dots, X_n)$ and take their average.
But there's a problem---we know what $f(\theta \mid X_1, \dots, X_n)$
is \emph{proportional} to, but we know nothing
about the exact probabilities.
How do we actually sample from this distribution?

\begin{remark}
    To see why sampling from such a distribution might be hard,
    consider the following challenge.
    \begin{quote}
        \textbf{Challenge.} Randomly shuffle a deck of $52$ cards
        so that the probability of any given ordering
        is proportional to $\sum_{n = 1}^{52} a_n \cdot \pi^n$,
        where the $n^{\text{th}}$ card has rank $a_n$.
    \end{quote}
    It's not computationally feasible to sum up all $52!$ weights.
\end{remark}

It turns out that there's a universal solution
to sampling from a distribution that we only know up to a proportion.
The solution is \textbf{Markov Chain Monte Carlo (MCMC)}.
We won't explain it here; for further details, see
\href{../../6.1220/notes-html/21.html}{Lecture 21 of 6.1220}.

\begin{quote}
    \textbf{Question \#2.}
    Determine how one might compute the maximum a posteriori
    $\hat{\theta}^{\text{MAP}} := \argmax_{\theta}
    f(\theta \mid X_1, \dots, X_n)$.
\end{quote}

The simplifying trick is to maximize the log of the posterior instead,
and do so via gradient ascent.  This looks like:
\[ \theta^{(k + 1)} = \theta^{(k)} +
\eta_k \cdot \left[
    \nabla \ell_n(\theta^{(k)}) + \nabla \log f(\theta^{(k)})
\right]. \]

\begin{quote}
    \textbf{Question \#3.}
    Determine how one might construct a \textbf{posterior interval}.
\end{quote}

Well, it's kinda obvious: a $(1 - \alpha)$ posterior interval
is any interval over which the integral of the posterior is $1 - \alpha$.

\subsection{Gaussian Conjugate Priors}

\begin{problem}
    Suppose our data is sampled i.i.d. as
    $X_1, \dots, X_n \sim \mathcal{N}(\theta, 1)$,
    and our prior is $\theta \sim \mathcal{N}(0, \sigma^2)$.
    Compute the posterior.
\end{problem}

\begin{solution}
    Plug it into the posterior formula.
    \begin{align*}
        f(\theta \mid X_1, \dots, X_n)
        & \propto \exp\left(
            -\frac{1}{2} \cdot \sum_{i = 1}^n (X_i - \theta)^2
        \right) \cdot \exp\left(
                -\frac{1}{2} \cdot \frac{\theta^2}{\sigma^2}
        \right) \\
        & = \exp\left(-\frac{1}{2} \left[
            \left[n + \frac{1}{\sigma^2}\right]\theta^2
            - 2\theta \sum_{i = 1}^n X_i
        \right]\right) \cdot
        \textcolor{teal}{
            \exp\left(-\frac{1}{2} \cdot \sum_{i = 1}^n X_i^2\right)
         } \\
        & \propto \exp\left(
            -\frac{\left[n + \frac{1}{\sigma^2}\right]}{2} \left(
                \theta - \frac{\sum_{i = 1}^n X_i}{
                    \left[n + \frac{1}{\sigma^2}\right]
                }
            \right)^2
        \right) \cdot
        \textcolor{teal}{\text{[constant, non-functions of $\theta$]}}
    \end{align*}
    We can throw away the constant, non-functions of $\theta$
    that accumulate when completing the square.  The result is:
    \[ \text{posterior} = \mathcal{N}\left(
        \frac{\sum_{i = 1}^n X_i}{\left[n + \frac{1}{\sigma^2}\right]},
        \frac{1}{\left[n + \frac{1}{\sigma^2}\right]}
    \right). \]
\end{solution}

There's not much insightful commentary to give here,
other than pointing out the obvious conjugate prior.

\end{document}